% =============================================================================
% Introduction and Related Work
% Working title: "Not All Judges Are Created Equal: Scoring Strategy
%                 No Universal Scoring Strategy for LLM-Judge Data Curation"
% =============================================================================

\section{Introduction}

Large language models are now routinely deployed as judges of other models'
outputs, and one of the highest-stakes uses of these judgments is training
data curation: scoring candidate instruction--response pairs and retaining
only those that exceed a quality bar before supervised fine-tuning. The
practice is ubiquitous---quality filtering with LLM judges underlies
instruction-tuning corpora~\citep{chen2024alpagasus, cui2023ultrafeedback},
multilingual pretraining pipelines~\citep{ali2025jql}, and synthetic data
generation at scale~\citep{sachdeva2024askllm}. Yet a basic design decision
in every such pipeline has gone almost entirely unexamined: \emph{how} the
judge is asked to express its judgment. Should it issue a binary
accept/reject verdict, a 1--5 Likert rating, a fine-grained 0--100 score, or
a comparative preference against a reference? Practitioners choose among
these formats by convention or convenience, implicitly assuming the choice
is innocuous. We show that it is not.

Prior work on judge scoring formats evaluates them against a single
criterion: agreement with humans or with other judges. Studies of grading
scales find that human--LLM alignment varies systematically with scale
granularity~\citep{li2025grading}, that absolute verbalized formats often
outperform comparative ones for system ranking~\citep{gera2025justrank}, and
that reference-anchored comparison reduces scoring variance across
judges~\citep{raee2026}. This evaluation-centric lens, however, measures the
wrong quantity for the curation setting. When judge scores drive data
selection, what matters is not whether the scores agree with human raters
but whether the \emph{selected subset} produces a better model after
fine-tuning. These two criteria can diverge: recent evidence shows that a
judge with moderate global correlation to reference labels can capture as
little as a fifth of the achievable gain in best-of-$n$ selection, because
selection depends on within-prompt discrimination while agreement metrics
are dominated by between-prompt variation, and because coarse scoring
formats produce ties on the majority of comparisons~\citep{landesberg2026}.
If agreement does not predict inference-time selection quality, there is
little reason to expect it to predict training-time curation quality, where
selection errors are compounded through gradient updates rather than
discarded after a single response.

We therefore pose the question the literature has skipped: \textbf{is
there a universal best scoring strategy for judge-driven data curation, or
does the downstream task dictate which strategy wins?} The field's
implicit assumption is universality---scale choices are made once,
justified (if at all) by human-alignment studies, and reused across tasks.
We hypothesize instead that the optimal strategy is
\emph{task-contingent}, for a mechanistic reason: scoring formats differ
in where they spend their discriminative capacity. Coarse formats (binary,
low-granularity Likert) act as high-precision detectors of categorical
failure but collapse gradations among acceptable responses; fine-grained
and comparative formats resolve those gradations at the cost of
calibration noise and verbosity-correlated drift. Tasks whose quality
structure is near-categorical---code that runs or does not, arithmetic
that is right or wrong---should therefore favor coarse verdicts, while
open-ended tasks with graded quality should favor formats that preserve
ordering among the acceptable. If this is right, no single scale can be
optimal everywhere, and the universality assumption silently degrades a
fraction of every curation pipeline built on it.

To test this, we treat the scoring format as the experimental treatment
and hold everything else fixed---the judge model, the candidate pool, the
selection budget, the fine-tuning recipe, and the evaluation protocol.
Five strategies span the design space used in practice: binary
accept/reject, 5-point Likert, 0--100 numeric, anchored comparison against
a fixed per-prompt reference, and within-prompt pairwise tournaments
aggregated with a Bradley--Terry model. For each strategy and retention
budget we select data by within-strategy percentile rank, fine-tune
identical base models across multiple seeds, and evaluate each resulting
model on multiple downstream tasks: judged instruction-following win-rates
(with an evaluation judge from a different model family than the curation
judge and position-swapped comparisons), and judge-free benchmarks of
constraint following and mathematical reasoning. Because ``no universal
best'' is a claim that cannot rest on a mere absence of a winner, our
pre-registered analysis requires positive evidence on two fronts:
significant within-task differences between strategies (Friedman tests
with Holm-corrected paired bootstraps), and instability of the strategy
ordering across tasks (top-1 winner agreement, Kendall's~$W$ concordance,
and a permutation test on the strategy$\times$task interaction).

Beyond the headline result, we analyze \emph{why} the winner moves.
Coarse formats induce massive score ties that force arbitrary tie-breaking
at the selection frontier; fine-grained formats spread scores but inject
calibration noise; comparative formats trade pointwise calibration for
within-prompt discrimination. We quantify these mechanisms through tie
probability, effective dynamic range, cross-strategy selection overlap,
and correlation with external reference annotations, and connect each
distributional property to the task structure that rewards or punishes it.

Our contributions are threefold. First, we provide, to our knowledge, the
first controlled comparison of judge scoring strategies evaluated by the
downstream performance of models fine-tuned on the data they select,
closing a loop that prior work leaves open at the agreement stage. Second,
we show that the optimal strategy is not universal: the strategy ranking
interacts with the downstream task, so scale choices optimized for human
alignment---or simply inherited by convention---can be actively suboptimal
for a given curation pipeline. From the mechanism analysis we distill
practical selection guidance mapping task properties (verifiability,
quality-distribution shape) to scoring formats. Third, we release an open
experimental harness\footnote{\url{https://github.com/lgondara/llm_judge_eval}}
with the pre-registered statistical machinery (paired bootstraps,
interaction permutation tests, concordance analysis) so the comparison
can be extended to new judges, domains, and formats at low cost.

% =============================================================================

\section{Related Work}

\paragraph{LLM-as-a-judge and scoring formats.}
Using LLMs to evaluate model outputs is now standard
practice~\citep{zheng2023mtbench, liu2023geval}, and a growing literature
examines how the elicitation format shapes the resulting
judgments. \citet{li2025grading} collect human and LLM ratings across 0--5,
0--10, and 0--100 scales on six benchmarks and find that human--LLM
alignment, measured by intraclass correlation, is highest on the coarsest
scale. \citet{gera2025justrank} compare scoring realizations---numeric,
Likert, comparative anchoring, and token probabilities---across 48 judges
and find that the realization affects system-level ranking quality nearly
as much as the choice of judge model itself, with verbalized absolute
formats generally outperforming comparative ones. Reference-anchored
schemes that convert judgments into win probabilities against a fixed
anchor reduce per-run standard error and cross-judge
variance~\citep{raee2026}, echoing earlier findings in human evaluation
that relative magnitude estimation yields far higher inter-annotator
reliability than Likert ratings~\citep{novikova2018rankme}. Related
analyses document scoring biases tied to the format
itself~\citep{li2025scoringbias} and show that operating on the judge's
full judgment distribution rather than a point estimate improves
inference~\citep{wang2025judgmentdist}. In information retrieval,
\citet{zhuang2024beyond} find that finer-grained relevance labels produce
better rankings than binary ones. All of these works share an
evaluation-centric criterion---agreement with humans, inter-judge
consistency, or ranking fidelity. None measures what happens when the
scores are consumed by a training pipeline.

\paragraph{Selection quality versus agreement.}
A small body of recent work questions whether agreement metrics predict
performance on the tasks judges are actually used
for. \citet{landesberg2026} show that on best-of-$n$ response selection a
judge with respectable global correlation recovers only a small fraction
of the oracle improvement, attributing the gap to weak within-prompt
discrimination and to ties induced by coarse pointwise scoring, and that
explicit pairwise judging recovers much of the lost signal.
\citet{balkir2026adaptive} develop adaptive testing machinery for
continuous judge scores, implicitly acknowledging that score granularity
interacts with ranking confidence. We extend this emerging
``selection $\neq$ agreement'' insight from inference-time reranking to
training-time data curation, where the consequences of selection error are
amplified: a misranked response is not merely served once but trained on,
shifting model parameters for every future query.

\paragraph{LLM-based data selection for instruction tuning.}
A parallel literature uses LLM judges to curate fine-tuning data but
treats the scoring format as fixed. AlpaGasus filters Alpaca with a
ChatGPT grader on a 0--5 scale and shows that 9k selected examples
outperform the full 52k~\citep{chen2024alpagasus}. UltraFeedback annotates
multi-model completions with GPT-4 scores to build preference
data~\citep{cui2023ultrafeedback}. DS\textsuperscript{2} models the error
structure of LLM ratings and corrects scores before selection, observing
that rating errors are prevalent and vary across rater
models~\citep{pang2025ds2}. Ask-LLM scores pretraining data with
instruction-tuned raters~\citep{sachdeva2024askllm}, and JQL extends
judge-based filtering to multilingual pretraining with demonstrated
downstream gains~\citep{ali2025jql}. Complementary LLM-free selectors rank
data by perplexity-derived difficulty~\citep{li2024ifd,
li2024superfiltering} or by diversity-aware
objectives~\citep{bukharin2024diversity}. Across this entire line of work
the judge's elicitation format is an unexamined constant: each paper
adopts one scale and compares against \emph{other selectors}, not against
the same selector under \emph{other scales}. The single closest precedent
we are aware of is an ablation showing four-bin relevance labels beat
two-bin labels for reranking in IR~\citep{zhuang2024beyond}---a result
confined to retrieval and never propagated through model training.

\paragraph{Positioning.}
Our work sits at the intersection of these threads and fills the gap
between them. From the scoring-format literature we inherit the treatment
variable; from the data-selection literature we inherit the task; from the
selection-versus-agreement literature we inherit the warning that the
standard evaluation criterion may be misleading. The resulting study is,
to our knowledge, the first to causally connect the choice of judge
scoring strategy to the quality of the fine-tuned model it ultimately
produces---and the first to test whether any such choice can be universal.
Agreement metrics are demoted from evaluation criterion to explanatory
mechanism, and the universality assumption embedded in every fixed-scale
curation pipeline is itself placed under test.

% =============================================================================
% BibTeX stubs (verify/complete before submission)
% =============================================================================
% @article{li2025grading,        title={Grading Scale Impact on LLM-as-a-Judge: Human-LLM Alignment Is Highest on 0-5 Grading Scale}, author={Li, Weiyue and others}, journal={arXiv:2601.03444}, year={2026}}
% @inproceedings{gera2025justrank, title={JuStRank: Benchmarking LLM Judges for System Ranking}, author={Gera, Ariel and others}, booktitle={ACL}, year={2025}}
% @misc{raee2026,                title={Reference-Anchored Elo Estimation}, note={OpenReview}, year={2026}}  % confirm publication status
% @inproceedings{novikova2018rankme, title={RankME: Reliable Human Ratings for Natural Language Generation}, author={Novikova, Jekaterina and Du{\v{s}}ek, Ond{\v{r}}ej and Rieser, Verena}, booktitle={NAACL}, year={2018}}
% @article{landesberg2026,       title={When LLM Judge Scores Look Good but Best-of-N Decisions Fail}, author={Landesberg, Eddie}, journal={arXiv:2603.12520}, year={2026}}
% @article{balkir2026adaptive,   title={Confident Rankings with Fewer Items: Adaptive LLM Evaluation with Continuous Scores}, author={Balk{\i}r, Esma and others}, journal={arXiv:2601.13885}, year={2026}}
% @inproceedings{chen2024alpagasus, title={AlpaGasus: Training a Better Alpaca with Fewer Data}, author={Chen, Lichang and others}, booktitle={ICLR}, year={2024}}
% @article{cui2023ultrafeedback, title={UltraFeedback: Boosting Language Models with High-Quality Feedback}, author={Cui, Ganqu and others}, journal={arXiv:2310.01377}, year={2023}}
% @inproceedings{pang2025ds2,    title={Improving Data Efficiency via Curating LLM-Driven Rating Systems}, booktitle={ICLR}, year={2025}}
% @article{sachdeva2024askllm,   title={How to Train Data-Efficient LLMs}, author={Sachdeva, Noveen and others}, journal={arXiv:2402.09668}, year={2024}}
% @article{ali2025jql,           title={JQL: Judging Quality across Languages}, year={2025}}  % verify
% @inproceedings{zheng2023mtbench, title={Judging LLM-as-a-Judge with MT-Bench and Chatbot Arena}, author={Zheng, Lianmin and others}, booktitle={NeurIPS}, year={2023}}
% @inproceedings{liu2023geval,   title={G-Eval: NLG Evaluation using GPT-4 with Better Human Alignment}, author={Liu, Yang and others}, booktitle={EMNLP}, year={2023}}
% @article{li2025scoringbias,    title={Evaluating Scoring Bias in LLM-as-a-Judge}, author={Li, Qingquan and others}, journal={arXiv:2506.22316}, year={2025}}
% @inproceedings{wang2025judgmentdist, title={Improving LLM-as-a-Judge Inference with the Judgment Distribution}, journal={arXiv:2503.03064}, year={2025}}
% @article{zhuang2024beyond,     title={Beyond Yes and No: Improving Zero-Shot LLM Rankers via Scoring Fine-Grained Relevance Labels}, author={Zhuang, Honglei and others}, year={2024}}
% @inproceedings{li2024ifd,      title={From Quantity to Quality: Boosting LLM Performance with Self-Guided Data Selection}, booktitle={NAACL}, year={2024}}
% @article{li2024superfiltering, title={Superfiltering: Weak-to-Strong Data Filtering for Fast Instruction-Tuning}, year={2024}}
% @article{bukharin2024diversity, title={Data Diversity Matters for Robust Instruction Tuning}, year={2024}}
